\section{预算会被花光}
\label{sec:17-Constrained-Guarantees/01-Differential-Privacy/03}

一个内部分析平台上线时做了件正确的事：所有查询都走差分隐私通道，每次按 \(\varepsilon=0.1\) 加噪，参数写在配置文件里，没人能绕过。上线公告写着``本平台满足 \(\varepsilon=0.1\) 的差分隐私''。

半年后审计。日志显示这半年跑了大约四千次查询，分布在几十个分析师手上，有的查总量，有的查分组，有的把同一个查询换了措辞重跑。审计员问了一个配置文件回答不了的问题：现在这个平台的隐私保证是多少？

答案是 \(400\)。这个数不是笔误——四千次 \(\varepsilon=0.1\) 的查询叠加，得到的是 \(\varepsilon=400\) 的保证。上一节讲过 \(\varepsilon\) 的语义是后验几率至多变化 \(e^{\varepsilon}\) 倍，\(e^{400}\) 这个数字大到没有任何含义。平台在数学上仍然``满足差分隐私''，实际上什么都没保护。

出问题的不是实现，是那句公告。\(\varepsilon=0.1\) 描述的是单次查询，而隐私不是一个可以挂在系统上的静态标签，它是一份\textbf{会被消耗的预算}。

\subsection{组合定理：账是加起来的}

\begin{theorem}[基本组合]
若机制 \(\mathcal M_1,\dots,\mathcal M_k\) 分别满足 \(\varepsilon_1,\dots,\varepsilon_k\)-差分隐私，则把它们依次作用在同一份数据上、并公开全部 \(k\) 个输出的复合机制满足 \(\big(\sum_i\varepsilon_i\big)\)-差分隐私。允许第 \(i\) 个机制的选择依赖前 \(i-1\) 个输出。
\end{theorem}

证明骨架：记完整输出为 \(t=(t_1,\dots,t_k)\)，按条件密度链式分解，
\[
\frac{p_D(t)}{p_{D'}(t)}
=\prod_{i=1}^{k}\frac{p_D(t_i\given t_1,\dots,t_{i-1})}{p_{D'}(t_i\given t_1,\dots,t_{i-1})}
\le\prod_{i=1}^{k}e^{\varepsilon_i}=e^{\sum_i\varepsilon_i}.
\]
每个因子的界来自第 \(i\) 个机制的差分隐私性质——注意此处已经把前面的输出固定住了，在给定它们的条件下 \(\mathcal M_i\) 是一个确定的机制，定义直接适用。这也是``允许自适应选择''这句话成立的原因：分析师看完第一个结果再决定第二个查什么，不改变每个条件因子的界。

条件对账：链式分解要求这 \(k\) 次都作用在\textbf{同一份数据}上；每个因子的界要求每个 \(\mathcal M_i\) 的保证对所有相邻对成立，而不只是对手上这份数据成立；\(\varepsilon_i\) 必须是事先固定的常数，如果分析师根据数据决定这次花多少预算，那个决定本身泄露信息，需要另行记账。

这条定理的形状很朴素，含义却很硬：\textbf{没有任何办法让重复查询免费}。有人会想，同一个计数查两次再取平均，噪声方差减半，精度提高——确实如此，但两次查询花掉 \(2\varepsilon\)，而直接用 \(2\varepsilon\) 单次查询能得到同样的精度。重复不产生额外的信息，只是把预算换成了精度，换算率还不占便宜。

这与\hyperref[sec:11-Mathematical-Statistics/05-Hypothesis-Testing/04]{``多重比较改变错误的计量单位''}是同一种结构。那里的问题是做了 \(m\) 次检验之后，``至少犯一次第一类错误''的概率远高于单次的 \(\alpha\)；这里的问题是做了 \(k\) 次查询之后，总泄露远高于单次的 \(\varepsilon\)。两处都在提醒：为单次操作设计的保证，不会因为使用者反复调用而自动延续。

\subsection{高级组合把线性改成平方根}

\(k\varepsilon\) 这个界是紧的——存在机制使它无法改进。但它紧在最坏情况上：\(k\) 次独立加噪，每次的密度比都恰好顶到 \(e^{\varepsilon}\) 且方向一致，这在概率上是极小事件。放弃对这个极小事件的控制，界就能大幅改善。

\begin{theorem}[高级组合]
对任意 \(\delta'>0\)，\(k\) 次 \(\varepsilon\)-差分隐私机制的自适应复合满足 \((\varepsilon',\delta')\)-差分隐私，其中
\[
\varepsilon'=\sqrt{2k\ln(1/\delta')}\;\varepsilon+k\varepsilon\,(e^{\varepsilon}-1).
\]
\end{theorem}

证明的核心是把对数密度比 \(\log\frac{p_D(t)}{p_{D'}(t)}\) 看成 \(k\) 个随机变量之和：每一项都被 \(\varepsilon\) 界住，期望是一个 \(O(\varepsilon^2)\) 的小量，于是这个和以极高概率集中在 \(\sqrt{k}\varepsilon\) 的量级上，只有指数小的概率跑到 \(k\varepsilon\) 附近。用的正是\hyperref[sec:05-Probability/07-Transforms-and-Bounds/04]{``Chernoff 方法与 Hoeffding 集中界''}里那套办法——把``最坏情况的和''换成``典型情况的和''，剩下的尾部交给 \(\delta'\) 吸收。

第二项 \(k\varepsilon(e^{\varepsilon}-1)\approx k\varepsilon^2\) 是那些对数密度比的期望累积起来的，\(\varepsilon\) 小的时候可以忽略，\(\varepsilon\) 接近 1 的时候会反过来主导，此时高级组合并不优于基本组合。这是套用公式时最常见的失误。

\begin{center}
\begin{tikzpicture}[x=0.038cm,y=0.55cm,>=stealth,font=\small]
  \draw[->,black!60] (-8,0) -- (268,0) node[right,font=\footnotesize] {查询次数 \(k\)};
  \draw[->,black!60] (0,-0.2) -- (0,7.3) node[above,font=\footnotesize] {累计 \(\varepsilon\)};
  \foreach \k in {50,100,150,200,250} {
    \draw[black!60] (\k,0) -- (\k,-0.16);
    \node[below,font=\footnotesize] at (\k,-0.18) {\k};
  }
  \draw[black!50,dashed] (0,5) -- (258,5);
  \node[anchor=west,font=\footnotesize,black!60] at (8,5.45) {预算上限};
  \draw[very thick,orange!85!black] (0,0) -- (70,7);
  \draw[very thick,blue!65] plot[domain=0:255,samples=80] (\x,{0.371*sqrt(\x)});
  \fill[black] (50,5) circle (2.2pt);
  \fill[black] (182,5) circle (2.2pt);
  \draw[black!55,dashed] (50,0) -- (50,5);
  \draw[black!55,dashed] (182,0) -- (182,5);
  \node[anchor=west,font=\footnotesize,orange!85!black] at (56,4.15) {\(k=50\) 撞线};
  \node[anchor=east,font=\footnotesize,blue!65] at (176,4.15) {\(k\approx180\) 撞线};
  \node[anchor=west,font=\footnotesize,orange!85!black] at (74,6.6) {基本组合 \(k\varepsilon\)};
  \node[anchor=west,font=\footnotesize,blue!65] at (120,2.6)
    {高级组合 \(\approx\sqrt{2k\ln(1/\delta')}\,\varepsilon\)};
\end{tikzpicture}
\end{center}

\emph{要记住的是两条曲线都在上升，只是速度不同。高级组合把撞线的位置从五十次推到一百八十次，并没有把上限推到无穷远——预算被花光只是时间问题，区别在于早晚。}

图里那条水平线的位置需要一个人来定，而这恰恰是最难的部分。\(\varepsilon=1\) 还是 \(\varepsilon=10\)，没有任何定理能给出答案，它取决于这份数据的敏感程度、发布对象是谁、被重识别之后的后果有多严重。已部署系统公布的总预算跨了两个数量级，各家的理由都不完全可比。这一步是政策判断而非数学判断，把它伪装成技术参数是不诚实的。

\subsection{后处理不变：一条很强、证明很短的性质}

\begin{proposition}[后处理不变性]
若 \(\mathcal M\) 满足 \((\varepsilon,\delta)\)-差分隐私，\(g\) 是任意（可随机的）映射且不再访问原始数据，则 \(g\circ\mathcal M\) 仍满足 \((\varepsilon,\delta)\)-差分隐私。
\end{proposition}

对确定性的 \(g\)，取任意输出集 \(S\)，令 \(T=g^{-1}(S)\)，那么 \(P(g(\mathcal M(D))\in S)=P(\mathcal M(D)\in T)\)，而定义对 \(T\) 这个集合已经成立。随机的 \(g\) 只需再对 \(g\) 的内部随机性取一次期望，不等式逐点保持。

一行就证完，结论却非常有用：加噪之后的数据可以随意分析、可视化、拿去训练模型、和其他公开数据交叉比对，隐私保证一分不减。这条性质是差分隐私区别于早期匿名化方案的关键——上一节那个交叉比对攻击之所以能成功，正是因为``删列''这类方案不具备后处理不变性，攻击者引入新的外部数据就能突破。

边界同样要说清。后处理不变的前提是 \(g\) \textbf{不再触碰原始数据}。工程实践里这条边界经常被无意越过：拿加噪结果去和真实数据做一致性校验，用真实的总数去归一化加噪的直方图，或者发现加噪结果``不合理''于是回去看一眼原始数据再调整——每一次回看都是一次新的数据访问，都要按组合定理记账。

组合的账还要分场合。\(k\) 次查询作用在同一份数据上是串行组合，\(\varepsilon\) 相加。如果把数据集切成互不相交的若干块，每块上各做一次 \(\varepsilon\)-差分隐私的查询，那是并行组合，总保证仍是 \(\varepsilon\) 而不是 \(k\varepsilon\)——因为改动一条记录只影响它所在的那一块，其余各块的输出分布完全没变。按地区分组统计属于并行组合，按不同疾病重复统计全体人群属于串行组合，两者形式相似而账目相差 \(k\) 倍。分不清这一点，要么白白浪费预算，要么严重高估自己的保证。

\subsection{固定预算下提升精度，只能靠结构}

预算既然会花光，自然要问能不能在同样的 \(\varepsilon\) 下把结果做得更准。答案是能，但改进的来源必须说清楚——所有有效的方法都在\textbf{利用问题本身的结构}，没有一个在绕过组合定理。

如果只关心一批查询里哪几个超过了阈值，稀疏向量技术让不超阈值的查询几乎不花预算，只有报告``超了''时才扣费；前提是真正超阈值的查询数量很少，这个稀疏性假设不成立时它不省任何东西。如果查询之间高度相关——比如一批嵌套的区间计数——可以先在一组精心构造的中间量上加噪，再线性组合出所有答案，利用的是这些查询张成的空间维数远低于查询个数。如果有可用的公开先验，可以用它给出一个不花预算的初始估计，只把预算用在修正上。

这些方法的共同点是：它们改善的是同一 \(\varepsilon\) 下的效用，而不是同一效用下的 \(\varepsilon\)。隐私与效用之间那条边界没有被推动，只是在既有边界上换了一个更合适的工作点。任何声称能同时改善两者而又不引入新假设的方案，要么用了未说明的结构假设，要么算错了账。

\subsection{DP-SGD：把训练循环接进这套记账}

训练一个模型可以看成几万次对数据的访问，组合定理直接适用，这就是 DP-SGD 的全部逻辑。

每一步先对\textbf{单个样本的梯度}做裁剪：\(g_i\leftarrow g_i\cdot\min\{1,\;C/\norm{g_i}_2\}\)。裁剪不是为了稳定训练，是为了把 \(\ell_2\) 敏感度\textbf{定义出来}——一条记录的进出至多改变梯度和 \(C\)，于是 \(\Delta_2=C\)。上一节讲过，没有有界的敏感度就没有可用的机制，所以这一步不可省略。然后按 Gaussian 机制给累加后的梯度加噪，除以批大小，交给优化器。最后按组合定理把每一步的 \((\varepsilon,\delta)\) 累计成总预算。

三处麻烦值得预先知道。裁剪阈值 \(C\) 取小了梯度方向被系统性扭曲，取大了噪声按 \(C\) 线性放大，形状与上一节那个截断边界的问题相同。逐样本梯度无法用常规的批量反向传播一次算出，实现上要么循环、要么用专门的技巧，训练时间通常翻几倍——这是真正的计算开销，与隐私记账无关。噪声注入的位置在梯度上，它与\hyperref[sec:14-Deep-Learning/03-Optimization-and-Training/02]{``动量与 Adam 在平均什么''}里那些依赖梯度二阶矩的自适应方法会互相干扰，因为二阶矩估计里现在混进了人为噪声。而在\hyperref[sec:14-Deep-Learning/03-Optimization-and-Training/03]{``归一化与 Dropout 改变训练时的随机系统''}的意义上，批归一化的批统计量跨样本耦合，会破坏``一条记录只影响自己那份梯度''的前提，DP 训练里通常换成层归一化或组归一化。

至于噪声对收敛的影响，\hyperref[sec:09-Convex-Optimization/08-Nonconvex-and-Stochastic-Optimization/04]{``SGD 在期望中下降但噪声会留下平台''}已经把机制讲过：注入的噪声不随步长衰减，最终会在最优点附近留下一个不收缩的平台，平台高度随 \(\sigma\) 上升。DP 训练的模型精度低于非私有训练，根源就在这里，不是调参没调好。

\begin{remark}
公开报告的 DP 训练结果里，同一个 \(\varepsilon\) 下的精度差异很大，这些差异来自裁剪策略、批大小、预训练权重是否公开可得等一系列选择，目前没有理论能预测哪种组合更好。选型层面的建议基本都是经验规律，不要当成定理来用；而记账层面的结论——敏感度、组合、后处理——是定理，可以放心依赖。
\end{remark}

\subsection{隐私是这一部分的第一个样本}

回到开头那个平台。它真正需要的不是更好的加噪实现，是一本账：谁查过什么、每次花了多少、总额还剩多少、花完之后系统该拒绝服务还是重新采集数据。差分隐私提供的是记账规则，不是一劳永逸的护身符。

这个形状会在这一部分反复出现。想要的两件好事——``发布有用的统计''与``单个人的参与不可辨认''——不能同时拉满，而它们之间的换算关系是一条可以证明的不等式，不是工程上的不够努力。下一个单元换一个场景：想要模型在干净数据上准确，又想要它在被人为扰动的数据上仍然准确。那里同样会遇到一条不可兼得的界，只是这一次，麻烦的来源不是记账，而是高维空间的几何。
